% Part: first-order-logic
% Chapter: model-theory
% Section: overspill

\documentclass[../../../include/open-logic-section]{subfiles}

\begin{document}

\olfileid[pt]{mod}{bas}{ove}
\olsection{Transbordamento}

\begin{thm}
\ollabel{overspill} Se um conjunto $\Gamma$ de sentenças tem modelos
finitos arbitrariamente grandes, então tem um modelo infinito.
\end{thm}

\begin{proof}
Amplie a linguagem de $\Gamma$ acrescentando uma quantidade enumerável
de novas constantes $c_0$, $c_1$, \dots e considere o conjunto $\Gamma
\cup \{c_i \neq c_j : i \neq j\}$. Dizer que $\Gamma$ tem modelos
finitos arbitrariamente grandes significa que, para todo $m >0$, existe
$n\ge m$ tal que $\Gamma$ tem um modelo de cardinalidade~$n$. Isso
implica que $\Gamma \cup \{c_i \neq c_j : i \neq j\}$ é finitamente
satisfazível. Pela compacidade, $\Gamma \cup \{c_i \neq c_j : i \neq
j\}$ tem um modelo $\Struct M$ cujo domínio deve ser infinito, pois ele
satisfaz todas as desigualdades $c_i \neq c_j$.
\end{proof}

\begin{prop}
\ollabel{inf-not-fo} 
Não existe sentença $!A$ de nenhuma linguagem de primeira ordem que seja
verdadeira em !!a{structure}~$\Struct M$ se, e somente se, o domínio
$\Domain{M}$ da !!{structure} for infinito.
\end{prop}

\begin{proof}
Se existisse tal $!A$, sua negação $\lnot !A$ seria verdadeira exatamente
nas !!{structure}s finitas e, portanto, teria modelos finitos
arbitrariamente grandes, mas não teria um modelo infinito, contradizendo
\olref{overspill}.
\end{proof}

\end{document}
